% !TeX root = height-and-moments-on-abelian0915.tex
\section{height-and-moments-on-abelian0915}

Plan:
\begin{enumerate}
	\item Heights (and Arakelov geometry);
	\item Berkovich space (and tropical space);
	\item Some relation between above two.
\end{enumerate}

\subsection[status=draft,tag=0000000,label=\labelNFM{subsec:height}]{Heights}

Given $a/b \in \bbQ$, $h(a/b) = $ ``number of bits needed to store $a/b$ on a computer''.
More explicitly, $h(a/b) \coloneqq \max \{ \log |a|, \log |b| \}$.

Let $K$ be a number field.
The height is
\[ h: \bbP^n(\bar{K}) \to [0,\infty), \quad [a_0:\cdots:a_n] \mapsto \sum_{v \in M_K} \frac{[K_v:\bbQ_v]}{[K:\bbQ]} \max_i \{ \log |a_i| \}. \]

\begin{theorem}[label=\labelNFM{thm:northcott}, status=draft, tag=0000000, title={Northcott property}]
	\[ \# \left\{ p \in \bbP^n(\bar{K}) \;\middle|\; \deg p \leq d, h(P) \leq B \right\} \]
	is finite.
\end{theorem}

\subsection[status=draft,tag=0000000,label=\labelNFM{subsec:heightmachine}]{Weil Height machine}

Let $X$ be a smooth projective variety over $K$.
Then there exists a unique homomorphism
\[ h_X: \Pic(X) \to \bbR^{X(\bar{K})}/\{O(1)\} \]
such that if $L$ is very ample with closed embedding $\phi_L: X \to \bbP^n$, $h_X(L) = h_{\bbP^n} \circ \phi_L + O(1)$.
And it is ``functorial'' and ``positive''.


\subsection[status=draft,tag=0000000,label=\labelNFM{subsec:nextsect}]{Canonical height}

Let $A$ be an abelian variety over $K$ and a line bundle on $A$.

\begin{theorem}[label=\labelNFM{thm:ntheight}, status=draft, tag=0000000]
	In the class of Weil heights, $h_L: A(\bar{K}) \to \bbR$ associated to $L$, there exist a unique $\hat{h}_L: A(\bar{K}) \to \bbR$ that can be written as a sum of quadratic and a linear form.
\end{theorem}

\begin{corollary}[label=\labelNFM{cor:c1}, status=draft, tag=0000000]
	On torsion points $x$, $\hat{h}_L(x) = 0$.
\end{corollary}


\subsection[status=draft,tag=0000000,label=\labelNFM{subsec:adelicintersection}]{Adelic intersection point of view}

Let $A$ be an abelian variety over $K$, $L$ a (rigidified) ample symmetric line bundel.
For each place, let $\|\cdot\|_{v,L}$ be the canonical metric, making $[m]^*L \cong L^{\otimes m^2}$ isometry on $L_v^\an$ over $A_v^\an$.
Set $\bar{L} = (L,(\|\cdot\|_{v,L})_{v \in M_K})$ be the adelic line bundle.

Let $Z$ be an integral cycle on $A$, (after Philippon, Guber, Bost-Gillet-Soule,Zhang), there exists an intersection theory.

\begin{definition}[label=\labelNFM{def:hi}, status=draft, tag=0000000]
	The canonical (Neron-Tate) height of $Z$ with respect to $L$ can be defined as
	\[ \hat{h}_L(Z) \coloneqq \frac{1}{[K:\bbQ]} \cdot \frac{1}{(\dim Z + 1) \deg_L Z} \left<\bar{L} \cdots \bar{L} | Z \right>. \]
	It is independent of the rigidified.
\end{definition}

\begin{example}[label=\labelNFM{eg:po}, status=draft, tag=0000000]
	If $Z$ is a point $P$, then we get again the canonical height.
\end{example}

\begin{example}[label=\labelNFM{eg:ag}, status=draft, tag=0000000]
	If you have a principally polarized abelian variety $(A,\lambda)$.
	That is, given a line bundle $L \cong \calO(\theta)$ with $\theta$ a divisor with one-dimensional global sections.
	Then
	\[ \hat{h}(\theta) \eqqcolon h_{NT}(\theta) = \frac{1}{[K:\bbQ]} \cdot \frac{1}{g \cdot g!} \left<\bar{L} \cdots \bar{L} | \theta \right> \]
\end{example}

\subsection[status=draft,tag=0000000,label=\labelNFM{subsec:Fali}]{Faltings height}

Let $A$ be an abelian variety over $K$.
($K$ is large enough)
We have
\[
	\begin{tikzcd}
		A & \calA \\
		\Spec K & \Spec \calO_K
		\arrow[from=1-1, to=1-2]
		\arrow[from=1-1, to=2-1]
		\arrow[from=1-2, to=2-2]
		\arrow[from=2-1, to=2-2]
	\end{tikzcd}
\]
It $(\calA)$ exists and called the Neron model.
And you have a zero section $0: \Spec \calO_K \to \calA$.
There is a line bundle
\[ \omega \coloneqq 0^* \Omega_{\calA/\Spec \calO_K}^{\wedge g} \]
on $\Spec \calO_K$.
It is equipped a canonical metric by Faltings.
Then the Faltings height is defined as
\[ h_F(A) = \frac{1}{[K:\bbQ]} \adeg \omega. \]

\begin{remark}[label=\labelNFM{rmk:111}, status=draft, tag=0000000]
	Using ``recursive formula'' of intersection numbers and Chambert-Loir measure, we have
	\[ \left< \bar{L} \cdots \bar{L} | Z \right> = \left< \bar{L} \cdots \bar{L} | \divisor_Z(s) \right> - \sum_{v \in M_K} \int_{A_v^\an} \log \|s\| \upd \mu. \]
	One obtains
	\[ \hat{h}_L(P) = \frac{1}{[K:\bbQ]} \sum_{v \in M_K} - \log \|s(P)\|_{v,L}, \quad P \notin \Supp D. \]
	And
	\[ h_{NT}(\theta) = \frac{1}{[K:\bbQ]g} \sum_{v\in M_K} \int_{A_v^\an} \log \|s\|_{v,L} \upd \mu.\]
\end{remark}


These are examples of the following principle:
\begin{slogan}
	Global quantities should be made from the local quantities.
\end{slogan}

\begin{theorem}[label=\labelNFM{thm:neron}, status=draft, tag=0000000]
	Let $v$ be a place of $K$, $A$ an abelian variety over $K$, $L$ a symmetric ample line bundle on $A$.
	There exists a unique up to constant $\Lambda: A(K_v)\setminus \Supp D \to \bbR$ such that
	\begin{enumerate}
		\item $\Lambda$ is continous;
		\item additive
		\item functorial w.r.t. isogeny.
		\item $\hat{h}_D(P) = \sum \Lambda_{D,v}(P) - C$ for $P \not\in \Supp D$.
	\end{enumerate}
\end{theorem}

% \begin{remark}[label=\labelNFM{rmk:rem}, status=draft, tag=0000000]
%   Consider
% \end{remark}


\subsection[status=draft,label=\labelNFM{subsec:ana}]{Analytic geometry over non-archimedean fields}

Let $(K,|\cdot|)$ be a non-archimedean fields, complete with respect to $|\cdot|$.
For example, $K = \bbQ_p$ or $\bbC_p$.
The topology of $K$ is totally disconnected.













































